\documentclass[12pt]{article}
\usepackage[T1]{fontenc}
\usepackage[utf8]{inputenc}
\usepackage{lmodern}
\usepackage{textcomp}
\usepackage{enumitem}
\usepackage{geometry}
\geometry{margin=1in}
\begin{document}
\begin{center}
Answer Key - Mathematics - Graduate Level Assessment 17 \
\small (Answers are ordered exactly as in the question paper.)
\end{center}

\section*{Multiple Choice}
\begin{enumerate}[label=\arabic*.]
\item \textbf{Answer:} B
\\ \textit{Options:} A) 100; B) 20; C) 14; D) 40; E) 80
\\ \textit{Note:} Correct option: B - 20
\item \textbf{Answer:} A
\\ \textit{Options:} A) 5 gallons; B) 6.2 gallons; C) 7.5 gallons; D) 7 gallons; E) 5.8 gallons
\\ \textit{Note:} Correct option: A - 5 gallons
\item \textbf{Answer:} A
\\ \textit{Options:} A) 4; B) 3; C) 1; D) 18; E) 51
\\ \textit{Note:} Correct option: A - 4
\item \textbf{Answer:} A
\\ \textit{Options:} A) [1, 2]; B) [0, 1]; C) [1, 1]; D) [4, 0]; E) [3, 0]
\\ \textit{Note:} Correct option: A - [1, 2]
\item \textbf{Answer:} A
\\ \textit{Options:} A) 4; B) 1; C) 9; D) 3; E) 2
\\ \textit{Note:} Correct option: A - 4
\end{enumerate}

\section*{True / False}
\begin{enumerate}[label=\arabic*.]
\setcounter{enumi}{5}
\item \textbf{Answer:} False
\\ \textit{Options:} A) True; B) False
\\ \textit{Note:} LLM-generated false statement. Correct answer: '15000'.
\item \textbf{Answer:} False
\\ \textit{Options:} A) True; B) False
\\ \textit{Note:} LLM-generated false statement. Correct answer: '$\frac{1}{4}$'.
\item \textbf{Answer:} False
\\ \textit{Options:} A) True; B) False
\\ \textit{Note:} LLM-generated false statement. Correct answer: '150'.
\item \textbf{Answer:} False
\\ \textit{Options:} A) True; B) False
\\ \textit{Note:} LLM-generated false statement. Correct answer: '-1.1'.
\item \textbf{Answer:} False
\\ \textit{Options:} A) True; B) False
\\ \textit{Note:} LLM-generated false statement. Correct answer: '-71'.
\end{enumerate}

\section*{Long Form Response}
\begin{enumerate}[label=\arabic*.]
\setcounter{enumi}{10}
\item \textbf{Answer:} We have $x<6+5=11$ and $x>6-5=1$, so the possible values of $x$ are from $2$ to $10$, and their difference is $10-2 = \boxed{8}$.
\\ \textit{Note:} Students should show all steps in their mathematical solution.
\item \textbf{Answer:} In order for a player to have an odd sum, he must have an odd number of odd tiles: that is, he can either have three odd tiles, or two even tiles and an odd tile. Thus, since there are $5$ odd tiles and $4$ even tiles, the only possibility is that one player gets $3$ odd tiles and the other two players get $2$ even tiles and $1$ odd tile. We count the number of ways this can happen. (We will count assuming that it matters in what order the people pick the tiles; the final answer is the same if we assume the opposite, that order doesn't matter.) $\dbinom{5}{3} = 10$ choices for the tiles that he gets. The other two odd tiles can be distributed to the other two players in $2$ ways, and the even tiles can be distributed between them in $\dbinom{4}{2} \cdot \dbinom{2}{2} = 6$ ways. This gives us a total of $10 \cdot 2 \cdot 6 = 120$ possibilities in which all three people get odd sums. In order to calculate the probability, we need to know the total number of possible distributions for the tiles. The first player needs three tiles which we can give him in $\dbinom{9}{3} = 84$ ways, and the second player needs three of the remaining six, which we can give him in $\dbinom{6}{3} = 20$ ways. Finally, the third player will simply take the remaining tiles in $1$ way. So, there are $\dbinom{9}{3} \cdot \dbinom{6}{3} \cdot 1 = 84 \cdot 20 = 1680$ ways total to distribute the tiles. We must multiply the probability by 3, since any of the 3 players can have the 3 odd tiles.Thus, the total probability is $\frac{360}{1680} = \frac{3}{14},$ so the answer is $3 + 14 = \boxed{17}$.
\\ \textit{Note:} Students should show all steps in their mathematical solution.
\end{enumerate}

\vfill
\noindent\textit{End of Answer Key}
\end{document}